\documentclass[conference]{IEEEtran}
\IEEEoverridecommandlockouts
\usepackage{cite}
\usepackage{amsmath,amssymb,amsfonts}
\usepackage{algorithm}
\usepackage{algorithmic}
\usepackage{graphicx}
\usepackage{textcomp}
\usepackage{xcolor}
\usepackage{booktabs}
\usepackage{multirow}
\usepackage{url}

\begin{document}

\title{Visual Grounding for Automated Reward Design:\\
When Does VLM Feedback Help LLM-Generated Reward Functions?}

\author{\IEEEauthorblockN{Anonymous Authors}}

\maketitle

% =============================================================
\begin{abstract}
Designing reward functions for reinforcement learning remains a major bottleneck
in robot manipulation. Recent work (Eureka, Ma et al.\ 2023) shows that large
language models (LLMs) can generate and iteratively refine reward code from task
descriptions and scalar training statistics. We extend this paradigm by closing
the loop with a vision-language model (VLM) that observes rollout videos of the
best candidate policy and provides behavioral diagnosis to the LLM before each
reward revision.

We evaluate the VLM+LLM method against an LLM-only (Eureka-style) baseline on
four ManiSkill manipulation tasks spanning a range of difficulty: PushCube
(easy), PickCube (easy), PushT (medium), and UnitreeG1PlaceAppleInBowl (hard,
6-stage humanoid bimanual). We also analyze supplementary VLM+LLM-only runs on
OpenCabinetDoor and OpenCabinetDrawer. Our main findings are:
(1)~On tasks that the LLM already solves at Iteration~0 (PushCube, PickCube),
VLM feedback is unnecessary and can even cause temporary regression by
prompting over-engineered reward designs.
(2)~On unsolved tasks (PushT, UnitreeG1), VLM feedback provides actionable
behavioral diagnosis---identifying specific failure modes such as
``press-and-freeze'' or stage-specific bottlenecks, and inferring behavioral
causes such as contact geometry and push direction from visual
observation---information that scalar metrics alone cannot reveal, leading to large performance gains (0\%$\to$56\% for PushT;
0\%$\to$87.5\% for UnitreeG1) compared to the LLM-only baseline (6.25\% and
12.5\%, respectively).
(3)~On partially solved cabinet tasks (Door/Drawer), VLM feedback provides an
early gain (both +12.5pp SE at Iter~1) but then plateaus: diagnoses remain
accurate yet repetitive, and challengers fail to surpass the elite.
(4)~The value of VLM feedback is phase-dependent: it is highest when the task is
unsolved or partially solved, and lowest when already solved. Here, ``progress''
means the current \emph{solution phase} measured by success\_at\_end (SE),
success\_once (SO), and re-evaluation stability, rather than PPO training step
count. Multi-stage tasks tend to provide qualitatively new VLM observations each
iteration, while simple single-stage tasks more often produce repetitive diagnoses.

These results suggest that VLM-in-the-loop reward design is most useful before
the policy is stably solved, and should be conditionally activated based on
current solution phase to avoid negative transfer on already-solved tasks.
\end{abstract}

\begin{IEEEkeywords}
reward design, vision-language models, reinforcement learning, manipulation,
LLM-guided optimization
\end{IEEEkeywords}

% =============================================================
\section{Introduction}

Reward engineering is one of the most labor-intensive aspects of applying
reinforcement learning (RL) to robotic manipulation. Hand-crafted reward
functions require domain expertise, extensive tuning, and often fail to
generalize across tasks. Recent advances in large language models (LLMs) have
opened a new avenue: using LLMs to \emph{automatically generate} reward
function code from natural-language task descriptions.

Eureka~\cite{eureka} demonstrated that an LLM can produce reward functions that
match or exceed human-designed ones across a range of tasks. The key insight is
an evolutionary outer loop: the LLM generates $K$ candidate reward functions per
iteration, each is trained with PPO in parallel, and the best candidate (by task
success rate) is selected. A ``reward reflection'' summary of training
statistics is fed back to the LLM for the next iteration.

However, a fundamental limitation of the Eureka approach is that the LLM
receives only \emph{scalar} feedback (success rate, return, per-component reward
statistics). When the task fails, the LLM knows \emph{that} it failed but not
\emph{how} or \emph{why}. For multi-stage manipulation tasks, this gap can be
critical: a scalar success rate of 0\% provides no signal about whether the
robot failed at reaching, grasping, transporting, or placing.

We propose closing this gap with a vision-language model (VLM) that observes
rollout videos of the best candidate's behavior and provides a structured
behavioral diagnosis. This diagnosis---identifying specific failure modes,
behavioral patterns, and suggested reward modifications---is injected into the
LLM's prompt alongside the scalar statistics before each reward revision.

Our contributions are:
\begin{enumerate}
  \item A VLM-in-the-loop extension of Eureka-style reward generation that adds
    visual behavioral feedback to the LLM's iterative refinement process
    (Section~\ref{sec:method}).
  \item An empirical evaluation on four manipulation tasks of varying difficulty,
    revealing that VLM feedback is most valuable on hard, multi-stage tasks
    and can be counterproductive on already-solved tasks
    (Section~\ref{sec:results}).
  \item Analysis of VLM feedback quality across solution phases,
    identifying failure modes of the VLM itself (severity miscalibration,
    repetitive diagnosis) and design guidelines for when to activate the visual
    feedback loop (Section~\ref{sec:analysis}).
\end{enumerate}


% =============================================================
\section{Related Work}

\textbf{LLM-based reward design.}
Eureka~\cite{eureka} uses GPT-4 to generate and iteratively refine reward
functions for RL, achieving human-level reward design on 29 tasks. L2R
(Language-to-Reward)~\cite{l2r} translates language instructions into reward
parameters. Text2Reward~\cite{text2reward} generates dense reward code from task
descriptions. Our work extends this line by adding a visual feedback channel.

\textbf{Vision-language models for robotics.}
VLMs have been used for task planning~\cite{saycan}, success
detection~\cite{vlm_success}, and policy evaluation. Recent work uses VLMs to
score trajectory quality or detect task completion. We use VLMs in a novel role:
providing behavioral \emph{diagnosis} that informs reward function design rather
than directly scoring or planning.

\textbf{Automated reward shaping.}
Classical approaches include potential-based reward shaping~\cite{pbrs} and
inverse RL. Our approach differs in that the reward function is generated as
executable code by an LLM, with the VLM providing qualitative behavioral
feedback rather than demonstrations.


% =============================================================
\section{Method}\label{sec:method}

\subsection{Overview}

Our method extends the Eureka outer loop with a VLM feedback channel. The
algorithm proceeds in $N$ iterations. Each iteration: (1)~the LLM generates $K$
candidate reward functions, (2)~each candidate is trained with PPO in parallel,
(3)~the best candidate is selected by task success rate, (4)~the VLM observes
a rollout of the best candidate and produces a behavioral diagnosis, and (5)~the
full history (code, metrics, VLM diagnosis) is fed back to the LLM.

\subsection{Algorithm}

\begin{algorithm}[t]
\caption{VLM-in-the-Loop Reward Generation}
\label{alg:vlm_loop}
\begin{algorithmic}[1]
\REQUIRE Task description $\mathcal{T}$, state docs $\mathcal{D}$, iters $N$, candidates $K$, RL steps $T$, LLM $\mathcal{L}$, VLM $\mathcal{V}$
\ENSURE Optimized reward function $R^*$
\STATE $\mathcal{H} \leftarrow \emptyset$, $R_{\text{best}} \leftarrow \texttt{sparse}$
\FOR{$n = 1$ \TO $N$}
    \STATE Build prompt $p_n$ from $\{\mathcal{T}, \mathcal{D}, R_{\text{best}}, \mathcal{H}\}$
    \STATE $\mathcal{C} \leftarrow \{R_1, \ldots, R_K\} \cup \{R_{\text{best}}\}$
    \COMMENT{$K$ new candidates + elite carry-over}
    \FOR{$R_k \in \mathcal{C}$ \textbf{in parallel}}
        \STATE Train $\pi_k$ with PPO for $T$ steps using $R_k$
        \STATE Evaluate $\pi_k$: $m_k = (\text{success\_once}, \text{success\_at\_end}, \text{return})$
    \ENDFOR
    \STATE $k^* \leftarrow \arg\max_k m_k.\text{success\_at\_end}$
    \STATE Record rollout of $\pi_{k^*}$, extract frames $\mathcal{F}$
    \STATE $c_{\text{vlm}} \leftarrow \mathcal{V}(\mathcal{F}, m_{k^*})$
    \COMMENT{Behavioral diagnosis}
    \STATE $\mathcal{H} \leftarrow \mathcal{H} \cup \{(n, R_{k^*}, m_{k^*}, c_{\text{vlm}}, \{m_k\}_k)\}$
\ENDFOR
\RETURN $R_{k^*}$
\end{algorithmic}
\end{algorithm}

\subsection{VLM Behavioral Diagnosis}

The VLM receives sampled frames from the best candidate's rollout along with
the episode's success label and scalar metrics. It is prompted to produce a
structured report with three sections: (1)~what the robot is currently doing,
(2)~what is going wrong (specific failure modes), and (3)~actionable reward
signal adjustments. This structured format ensures the diagnosis is directly
usable by the LLM for reward code modification.

\subsection{Elite Carry-Over}

Following Eureka, we carry the previous best candidate as an ``elite'' into the
next iteration's candidate pool. This provides a safety net: even if all new
LLM-generated candidates perform poorly, the best-so-far is preserved. Our
analysis shows this mechanism is critical when VLM feedback leads the LLM to
over-engineer solutions on easy tasks.

\subsection{Comparison: VLM+LLM vs Eureka (LLM-only)}

The Eureka baseline uses the same outer loop but without VLM feedback. The LLM
receives only scalar metrics and reward reflection statistics. Both methods use
identical hyperparameters ($K=4$, $N=5$, same PPO configuration and seed).


% =============================================================
\section{Experimental Setup}

\subsection{Tasks}

We evaluate on four ManiSkill~\cite{maniskill} manipulation tasks:

\begin{table}[htbp]
\caption{Task characteristics}
\label{tab:tasks}
\centering
\begin{tabular}{lccl}
\toprule
Task & Stages & Robot & Difficulty \\
\midrule
PushCube-v1 & 1 & Panda & Easy \\
PickCube-v1 & 2 & Panda & Easy \\
PushT-v1 & 1 & Panda & Medium \\
UnitreeG1PlaceApple & 6 & G1 humanoid & Hard \\
\bottomrule
\end{tabular}
\end{table}

\subsection{Supplementary Cabinet Runs}

To stress-test the phase-dependent hypothesis on articulated but largely
single-phase tasks, we additionally analyze OpenCabinetDoor-v1 and
OpenCabinetDrawer-v1 (VLM+LLM only; no clean Eureka counterparts available).
These runs are used for hypothesis refinement rather than primary
VLM-vs-Eureka benchmarking.

\subsection{Hyperparameters}

All experiments use $K=4$ candidates per iteration, $N=5$ outer iterations, and
$N_{\text{eval}}=16$ evaluation episodes per candidate. Fitness is
$\text{success\_at\_end}$ (fraction of episodes where the task is completed at
the final timestep). The LLM is GPT-4o; the VLM is also GPT-4o with vision.
Seed is fixed at 9351 for all runs.

% TODO: PPO hyperparameters table


% =============================================================
\section{Results}\label{sec:results}

\subsection{Main Results}

\begin{table}[htbp]
\caption{Best success\_at\_end (\%) across methods}
\label{tab:main}
\centering
\begin{tabular}{lccc}
\toprule
Task & VLM+LLM & Eureka & Improvement \\
\midrule
PushCube-v1 & 93.75 & --- & N/A \\
PickCube-v1 & 100.0 & --- & N/A \\
PushT-v1 & \textbf{56.25} & 6.25 & 9$\times$ \\
UnitreeG1PlaceApple & \textbf{87.50} & 12.50 & 7$\times$ \\
\bottomrule
\end{tabular}
\end{table}

% NOTE: PushCube and PickCube have no Eureka baseline runs yet.
% The reported VLM+LLM numbers are from the latest runs (2026-03-04).
% PushT and UnitreeG1 Eureka baselines had a VLM leak bug in Iter 0
% (stale VLM comment from the resumed VLM+LLM run). This has been fixed
% in code but the Eureka runs have not been re-executed yet.

On the two hard tasks, VLM+LLM achieves 7--9$\times$ higher success than the
Eureka baseline. On easy tasks, VLM+LLM achieves near-perfect performance from
Iteration~0, making VLM feedback unnecessary.

In supplementary cabinet runs, VLM feedback produces early improvements
(Door: 75.0\%$\to$87.5\%; Drawer: 62.5\%$\to$75.0\% at Iter~1), followed by
plateau. In Drawer, candidate timeouts reduced effective search width to $K=2$
throughout; in Door, candidate attrition reduced effective $K$ in Iter~2--3.

\subsection{Per-Iteration Trajectories}

% TODO: Figure showing SE over iterations for all 4 tasks, VLM+LLM vs Eureka

\textbf{UnitreeG1PlaceAppleInBowl} (hardest task): VLM+LLM progresses from 0\%
(Iter~0) to 100\% success\_once at Iter~2 and 87.5\% success\_at\_end at
Iter~3. Each VLM observation identifies a new stage-specific bottleneck:
``grasps but fails to transport'' (Iter~0), ``nudge-and-abandon'' regression
(Iter~1), ``near-bowl freeze'' (Iter~2--3). The Eureka baseline plateaus at
6.25--12.5\%.

\textbf{PushT-v1}: VLM+LLM reaches 56.25\% success\_once at Iter~4 after VLM
identifies a ``press-and-freeze'' failure pattern. Eureka achieves only 6.25\%.

\textbf{PushCube-v1}: Achieves 100\% SE at Iter~0. Subsequent VLM feedback
identifies ``overshoot / over-pushing'' every iteration, prompting complex
geometric strategies that perform worse. Best SE after 5 iterations: 93.75\%.
The apparent regression from 100\% may not be VLM-induced: the same Iter~0
policy re-evaluated at only 43.75\% in a later iteration, indicating that the
initial 100\% was likely inflated by N=16 evaluation noise rather than
representing true performance.

\textbf{PickCube-v1}: Maintains 93.75--100\% SE across all iterations. VLM
correctly reports ``no clear failure'' in 3/5 iterations, overreacts once
(``near-failure / brittle strategy'' on a 100\% SE episode). The final best
candidate's anti-stall code originated from VLM's one useful diagnosis.


% =============================================================
\section{Analysis}\label{sec:analysis}

\subsection{When Does VLM Feedback Help?}

\begin{table}[htbp]
\caption{VLM effectiveness by task and solution phase}
\label{tab:vlm_effectiveness}
\centering
\begin{tabular}{lccl}
\toprule
Task & Iter 0 SE & VLM Effect & Diagnosis Type \\
\midrule
PushCube & 100\% & Regression$^*$ & Stylistic (non-blocking) \\
PickCube & 100\% & Neutral & Mostly ``no failure'' \\
PushT & 0\% & Improvement & Blocking failure mode \\
UnitreeG1 & 0\% & Major gain & Stage bottleneck \\
OpenCabinetDoor$^\dagger$ & 75\% & Early gain, then plateau & Partial-open stall / contact drift \\
OpenCabinetDrawer$^\dagger$ & 62.5\% & Early gain, then plateau$^\ddagger$ & Partial-open retreat / disengagement \\
\bottomrule
\end{tabular}
\end{table}

\noindent{\footnotesize $^*$Observed regression is confounded with N=16 evaluation noise (same policy re-evaluated $\pm$56pp).}
\noindent{\footnotesize $^\dagger$Supplementary VLM+LLM-only analysis (no clean Eureka counterpart).}
\noindent{\footnotesize $^\ddagger$Drawer had effective $K=2$ due repeated candidate timeouts; Door also had candidate attrition in late iterations.}

The value of VLM feedback is better explained by the policy's \emph{solution
phase} than by task label alone. We define progress as:
\begin{itemize}
  \item \textbf{Unsolved phase:} low SE and low SO (policy rarely reaches the goal).
  \item \textbf{Partially solved phase:} SO high but SE clearly lower (policy reaches but does not reliably hold/finish).
  \item \textbf{Solved phase:} high SE with stable re-evaluation across iterations.
\end{itemize}
Under this definition, VLM feedback is typically most useful in unsolved and
partially solved phases, and often less useful (or harmful) in the solved phase.
The cabinet runs refine this claim: in partially solved settings, VLM can still
improve performance early, but benefits may saturate quickly when failure
diagnoses become repetitive and search width is reduced by invalid candidates.

\subsection{VLM Failure Modes}

We identify three failure modes of the VLM feedback loop:

\textbf{Severity miscalibration.} The VLM is prompted to find ``failure modes''
and produces diagnoses with similar urgency regardless of success rate. On
PushCube (100\% SE), it reports ``overshoot / lack of termination'' with the
same alarm as on UnitreeG1 (0\% SE). This causes the LLM to treat non-blocking
stylistic issues as critical failures requiring architectural reward changes.

\textbf{Repetitive diagnosis on simple tasks.} PushCube's VLM produces the same
``overshoot'' observation across all 5 iterations. Unlike UnitreeG1 where VLM
discovers qualitatively new bottlenecks each iteration (grasp $\to$ transport
$\to$ place $\to$ release), simple tasks have only one observable behavior
pattern.

\textbf{Over-engineering induction.} VLM diagnoses on successful tasks prompt
the LLM to generate complex reward strategies (explicit geometric staging,
curriculum gates, behind-pose requirements) that perform worse than the simple
potential-based reward that already works. On PushCube, every iteration produced
at least one 0\% SE candidate inspired by VLM-suggested complexity.

\textbf{Search-width confound in plateau regimes.} In cabinet runs, multiple
candidates failed to complete training (timeouts/runtime failures), reducing
effective $K$ and lowering the chance of finding improvements in later
iterations. Thus, late-iteration plateau should be interpreted as a joint effect
of diagnosis saturation and reduced search width, not solely as a VLM failure.

\subsection{Reward-Task Misalignment vs Reward Hacking}

A recurring pattern across tasks is high return with low success. We note that
this is better described as \emph{reward-task misalignment} rather than
``reward hacking'': the agent rationally optimizes the reward function, but the
reward fails to adequately capture the task objective. The agent is not
exploiting loopholes; the reward is simply misspecified. VLM feedback helps
detect and diagnose such misalignment on hard tasks by observing the specific
behavior the agent adopts under the misaligned reward.

\subsection{N=16 Evaluation Noise}

With only 16 evaluation episodes, success rate measurements have a granularity
of 6.25\% and substantial stochastic variance. We observe elite carry-over
candidates whose SE fluctuates by up to 56 percentage points across
re-evaluations of the identical policy (PushCube: 100\%$\to$43.75\%). This
noise affects both methods equally but complicates iteration-level comparisons.
% TODO: Consider increasing to N=64 or N=100 for final experiments.

\subsection{Role of Elite Carry-Over}

The elite mechanism acts as a critical safety net against VLM-induced regression
on easy tasks. Without it, the best PushCube solution (100\% SE at Iter~0) would
have been permanently lost when VLM-guided candidates underperformed. In cabinet
runs, elite carry-over preserved the Iter~1 breakthroughs (Door 87.5\%, Drawer
75\%) during post-breakthrough plateau. On hard tasks, the elite preserves gains
between breakthrough iterations.


% =============================================================
\section{Discussion and Limitations}

\textbf{When to use VLM feedback.}
Our results suggest phase-aware activation of VLM feedback. In unsolved and
partially solved phases (low SE, or high SO but low SE), VLM diagnosis can
improve intervention targeting. In solved phases (high and stable SE), VLM
feedback more often emphasizes non-blocking issues and can induce
over-engineering. An adaptive system could gate VLM usage by phase and
terminate the VLM loop when diagnoses repeat without SE improvement.

\textbf{Single-seed evaluation.}
All experiments use a single seed (9351). While K=4 candidates per iteration
provide diversity in reward function design (the primary experimental variable),
the PPO training stochasticity is not independently sampled across seeds. We
argue this is acceptable because the reward function diversity---not seed
diversity---drives the key comparisons, and N=16 evaluation episodes provide
some statistical coverage. However, larger-scale multi-seed validation is
needed.

% TODO: Eureka baseline VLM leak bug. The current Eureka runs for PushT and
% UnitreeG1 have a stale VLM comment from Iter 0 leaked into the LLM prompt
% (from resuming a VLM+LLM run). The bug has been fixed in code but the
% baselines need to be re-run for a clean comparison. Current numbers should
% be interpreted as "fresh iterative VLM" vs "stale/frozen VLM" rather than
% a perfect VLM vs no-VLM ablation.

\textbf{VLM prompt design.}
The current VLM prompt always asks for ``failure analysis,'' which biases
it toward finding problems even on successful episodes. A severity-aware
prompt that conditions on the current success rate could reduce
over-engineering on easy tasks.

\textbf{Evaluation granularity.}
N=16 episodes is insufficient for fine-grained performance comparison.
The observed elite re-evaluation noise ($\pm$25 percentage points) limits
our ability to detect small improvements or confirm regressions.


% =============================================================
\section{Conclusion}

We presented a VLM-in-the-loop extension of Eureka-style automated reward
design. By closing the loop between policy behavior observation (VLM) and reward
code generation (LLM), we enable the system to diagnose and fix failure modes
that scalar training statistics alone cannot reveal.

Our main finding is that \textbf{VLM feedback value is phase-dependent}:
\begin{itemize}
  \item On hard, multi-stage tasks (UnitreeG1: 6 stages), VLM provides
    stage-specific bottleneck identification that drives 7$\times$ performance
    improvement over the LLM-only baseline (87.5\% vs 12.5\%).
  \item On medium tasks (PushT), VLM identifies behavioral failure modes
    invisible to scalar metrics, enabling 9$\times$ improvement (56.25\% vs
    6.25\%).
  \item On partially solved articulated tasks (OpenCabinetDoor/Drawer), VLM
    provides early gains (+12.5pp at Iter~1) but then plateaus; repetitive
    diagnoses and reduced effective search width limit further improvement.
  \item On easy, already-solved tasks (PushCube, PickCube), VLM feedback is
    unnecessary and may cause regression through over-engineering induction,
    though observed regressions are confounded with N=16 evaluation noise.
\end{itemize}

These results point to a design principle for VLM-in-the-loop systems:
\textbf{visual feedback should be conditionally activated by solution phase}
(unsolved / partially solved / solved), with early stopping when the VLM has no
new observations or no SE improvement.
The VLM excels at answering ``how is the robot failing?''---a question that only
matters when the robot \emph{is} failing.

% Future work: (1) VLM_full_RewardPlot extension combining rollout images with
% reward component plots for richer VLM diagnosis; (2) severity-aware VLM
% prompting; (3) multi-seed validation; (4) additional tasks (PegInsertion,
% OpenCabinetDoor).


% =============================================================
\begin{thebibliography}{00}
\bibitem{eureka} Y.~Ma, W.~Liang, G.~Wang, D.-A.~Huang, O.~Bastani,
  D.~Jayaraman, Y.~Zhu, L.~Fan, and A.~Anandkumar, ``Eureka: Human-level
  reward design via coding large language models,'' in \emph{Proc. ICLR}, 2024.
\bibitem{l2r} W.~Yu, N.~Gileadi, C.~Fu, S.~Kirmani, K.-H.~Lee, M.~G.~Arenas,
  H.-T.~L.~Chiang, T.~Erez, L.~Hasenclever, J.~Humplik, et al.,
  ``Language to rewards for robotic skill synthesis,'' \emph{arXiv:2306.08647}, 2023.
\bibitem{text2reward} T.~Xie, S.~Zhao, C.~H.~Wu, Y.~Liu, Q.~Luo, V.~Zhong,
  Y.~Yang, and T.~Yu, ``Text2reward: Reward shaping with language models for
  reinforcement learning,'' in \emph{Proc. ICLR}, 2024.
\bibitem{maniskill} J.~Gu, F.~Xiang, X.~Li, Z.~Ling, X.~Liu, T.~Mu,
  Y.~Tang, S.~Tao, X.~Wei, Y.~Yao, et al., ``ManiSkill2: A unified benchmark
  for generalizable manipulation skills,'' in \emph{Proc. ICLR}, 2023.
\bibitem{saycan} M.~Ahn, A.~Brohan, N.~Brown, Y.~Chebotar, O.~Cortes,
  B.~David, C.~Finn, C.~Fu, K.~Gopalakrishnan, K.~Hausman, et al., ``Do as I
  can, not as I say: Grounding language in robotic affordances,''
  \emph{arXiv:2204.01691}, 2022.
\bibitem{vlm_success} W.~Du, Y.~Chen, W.~L.~Chao, and D.~Jayaraman,
  ``Vision-language models as success detectors,'' \emph{arXiv:2303.07280}, 2023.
\bibitem{pbrs} A.~Y.~Ng, D.~Harada, and S.~Russell, ``Policy invariance under
  reward transformations: Theory and application to reward shaping,'' in
  \emph{Proc. ICML}, 1999.
\end{thebibliography}

\end{document}
